\documentclass[11pt,compress,t,notes=noshow, xcolor=table]{beamer}

\input{../../style/preamble}
\input{../../latex-math/basic-math}
\input{../../latex-math/basic-ml}
\input{../../latex-math/optim-basics}

\newcommand{\alphav}{\bm{\alpha}}
\newcommand{\betav}{\bm{\beta}}

\title{Optimization in Machine Learning}

\begin{document}

\titlemeta{
Constrained Optimization
}{
Lagrangian Relaxation and Duality
}{
figure/duality_minimax.pdf
}{
\item Relax constraints with Lagrange multipliers
\item Dual function, weak duality, and strong duality
\item Min-max ordering, saddle points, and KKT conditions
\item Example: deriving the hinge-loss SVM dual
}

\begin{framev}[fs=footnotesize]{Why Move to the Dual?}
\splitV[0.48]{
\textbf{Primal problem}
$$
\begin{aligned}
\min_{\xv \in \R^d} \quad & F(\xv) \\
\text{s.t.} \quad & g_i(\xv) \le 0, \quad i = 1,\dots,m
\end{aligned}
$$
\begin{itemize}
\item Primal methods enforce feasibility directly in the original variables.
\item Dual methods \textbf{relax} the constraints and attach nonnegative multipliers to violations.
\item Each fixed multiplier vector $\alphav$ yields a \textbf{lower bound}; maximizing over $\alphav$ tightens it.
\end{itemize}
}{
\imageC[0.67]{figure/duality_lower_bound.pdf}
}
\end{framev}


\begin{framei}[fs=small]{Lagrangian Relaxation}
\item For inequality constraints, define the Lagrangian minimax objective
$$
H(\xv,\alphav) := F(\xv) + \sum_{i=1}^m \alpha_i g_i(\xv),
\qquad
\alphav \ge \zero.
$$
\item The corresponding \textbf{dual function} is the relaxed optimum at fixed multipliers:
$$
L(\alphav) := \min_{\xv} H(\xv,\alphav).
$$
\item The minimization is only with respect to $\xv$; the multiplier vector $\alphav$ is held fixed.
\item Why $\alpha_i \ge 0$? If $g_i(\xv) > 0$, then the violated constraint increases the objective by $\alpha_i g_i(\xv)$.
\item If $g_i(\xv) < 0$, the term can decrease the objective, which is why $L(\alphav)$ is a lower bound rather than a direct replacement for the primal.
\item Equality constraints $h_j(\xv)=0$ are handled by adding $\sum_j \beta_j h_j(\xv)$ with sign-unrestricted multipliers $\beta_j \in \R$.
\end{framei}


\begin{framei}[fs=small]{Weak Duality}
\item For any feasible $\xv$ and any multiplier vector $\alphav \ge \zero$,
$$
H(\xv,\alphav)
=
F(\xv) + \sum_{i=1}^m \alpha_i g_i(\xv)
\le
F(\xv),
$$
because each $g_i(\xv) \le 0$ and each $\alpha_i \ge 0$.
\item Taking the minimum over $\xv$ on the left yields
$$
L(\alphav) = \min_{\xv} H(\xv,\alphav) \le F(\xv)
\qquad \text{for every feasible } \xv.
$$
\item In particular, if $p^\ast$ denotes the optimal primal value, then
$$
L(\alphav) \le p^\ast
\qquad \forall \alphav \ge \zero.
$$
\item Maximizing this lower bound gives the \textbf{dual problem}
$$
d^\ast
:=
\max_{\alphav \ge \zero} L(\alphav)
=
\max_{\alphav \ge \zero} \min_{\xv} H(\xv,\alphav),
$$
and therefore
$$
d^\ast \le p^\ast.
$$
\item This universal lower-bound property is called \textbf{weak duality}.
\end{framei}


\begin{framev}[fs=scriptsize]{Ordering Matters: Min-Max vs.\ Max-Min}
\textbf{Dual:}
$$
d^\ast = \max_{\alphav \ge \zero} \min_{\xv} H(\xv,\alphav)
$$
\textbf{Primal:}
$$
p^\ast = \min_{\xv} \max_{\alphav \ge \zero} H(\xv,\alphav)
$$
\begin{itemize}
\item Always, $\max_{\alphav} \min_{\xv} H(\xv,\alphav) \le \min_{\xv} \max_{\alphav} H(\xv,\alphav)$.
\item Without convex-concave structure, the order can change the value.
\item If $H$ is convex in $\xv$ and concave in $\alphav$, a saddle point makes the two orders equal.
\end{itemize}
\imageC[0.84]{figure/duality_minimax.pdf}
\end{framev}


\begin{framei}[fs=small]{The Primal Is the Reversed Minimax}
\item Fix the primal variable $\xv$ and maximize $H(\xv,\alphav)$ over $\alphav \ge \zero$.
\item If $\xv$ violates any constraint, say $g_k(\xv) > 0$, then
$$
\max_{\alphav \ge \zero} H(\xv,\alphav) = \infty,
$$
because we can send $\alpha_k \rightarrow \infty$.
\item If $\xv$ is feasible, then each penalty term satisfies $\alpha_i g_i(\xv) \le 0$, so the maximum is achieved at $\alphav = \zero$:
$$
\max_{\alphav \ge \zero} H(\xv,\alphav) = F(\xv).
$$
\item Therefore,
$$
\min_{\xv} \max_{\alphav \ge \zero} H(\xv,\alphav)
=
\min_{\xv \in \S} F(\xv)
=
p^\ast.
$$
\item Primal and dual thus optimize the \textbf{same} Lagrangian objective, but in different orders:
$$
p^\ast = \min_{\xv} \max_{\alphav \ge \zero} H(\xv,\alphav),
\qquad
d^\ast = \max_{\alphav \ge \zero} \min_{\xv} H(\xv,\alphav).
$$
\end{framei}


\begin{framei}[fs=small]{Strong Duality and Slater's Condition}
\item Weak duality always gives
$$
d^\ast \le p^\ast.
$$
\item To obtain equality, we usually need convex structure:
\begin{itemize}
\item $F(\xv)$ convex
\item each $g_i(\xv)$ convex
\item hence $H(\xv,\alphav)$ convex in $\xv$ and linear, therefore concave, in $\alphav$
\end{itemize}
\item A standard qualification is \textbf{Slater's condition}: there exists at least one strictly feasible point $\xv_0$ such that
$$
g_i(\xv_0) < 0
\qquad \forall i=1,\dots,m.
$$
\item Under convexity plus Slater's condition, the duality gap disappears:
$$
d^\ast = p^\ast.
$$
\item This is \textbf{strong duality}, and it is what makes dual optimization so useful in many machine-learning problems.
\end{framei}


\begin{framei}[fs=footnotesize]{KKT Conditions}
\item For differentiable convex problems under suitable regularity conditions, optimality is characterized by the \textbf{Karush-Kuhn-Tucker (KKT)} conditions.
\item For inequality-constrained problems, a pair $(\xv^\ast,\alphav^\ast)$ is optimal if and only if:
\item \textbf{Primal feasibility:}
$$
g_i(\xv^\ast) \le 0
\qquad \forall i
$$
\item \textbf{Dual feasibility:}
$$
\alphav^\ast \ge \zero
$$
\item \textbf{Complementary slackness:}
$$
\alpha_i^\ast g_i(\xv^\ast) = 0
\qquad \forall i
$$
\item \textbf{Stationarity:}
$$
\nabla F(\xv^\ast) + \sum_{i=1}^m \alpha_i^\ast \nabla g_i(\xv^\ast) = \zero
$$
\item Interpretation: inactive constraints have $\alpha_i^\ast = 0$, while active constraints can have nonzero multipliers.
\item The multipliers act like \textbf{shadow prices}: they measure how strongly the objective pushes against each active constraint.
\end{framei}


\begin{framev}[fs=scriptsize]{Toy Example: One Active Multiplier}
Reuse
\[
\min_{\wv}
F(\wv)
=
(w_1 - 1.2)^2
+
(w_1 - 1.2)(w_2 - 0.2)
+
(w_2 - 0.2)^2
\]
\[
\text{s.t. }
\qquad
0 \le w_1 \le 1,
\qquad
0 \le w_2 \le 0.9.
\]
Write the four box constraints as
\[
g_1(\wv)=-w_1,\;
g_2(\wv)=w_1-1,\;
g_3(\wv)=-w_2,\;
g_4(\wv)=w_2-0.9.
\]
\splitV[0.48]{
At the optimum $\wv^\ast=(1,0.3)^\top$,
\[
\nabla F(\wv^\ast)=(-0.3,0)^\top.
\]
Hence the KKT multipliers are
\[
\alphav^\ast=(0,0.3,0,0)^\top:
\]
only the right wall is active, and its multiplier balances the push of the objective.
}{
\imageC[1]{figure/duality_box_kkt.pdf}
}
\end{framev}


\begin{framei}[fs=small]{General Recipe for Deriving a Dual}
\item \textbf{1. Start from a constrained primal.} Introduce slack variables if they simplify the objective or the constraints.
\item \textbf{2. Form the Lagrangian} $H(\xv,\alphav,\betav)$ by attaching a multiplier to each constraint.
\item \textbf{3. Eliminate primal variables} by computing
$$
L(\alphav,\betav) = \min_{\xv} H(\xv,\alphav,\betav).
$$
Usually this is done via stationarity conditions with respect to the primal variables.
\item \textbf{4. Substitute back} to obtain a pure optimization problem in the dual variables, together with any remaining sign or box constraints.
\item \textbf{5. Solve the dual} and then recover the primal optimum from the stationarity relations.
\item In many important examples, deriving the dual is mostly algebra: differentiate, substitute, simplify, and interpret the remaining constraints.
\end{framei}


\begin{framei}[fs=footnotesize]{ML Example: SVM Primal and Stationarity}
\item Start from the soft-margin SVM with slack variables:
$$
\min_{\wv,b,\xi_1,\dots,\xi_n}
\frac{1}{2}\|\wv\|_2^2 + C \sum_{i=1}^n \xi_i
$$
$$
\text{s.t. }
\qquad
1 - y^{(i)} \big({\xv^{(i)}}^\top \wv + b\big) - \xi_i \le 0,
\qquad
\xi_i \ge 0.
$$
\item Introduce multipliers $\alpha_i \ge 0$ for the margin constraints and $\gamma_i \ge 0$ for the nonnegativity constraints.
\item Stationarity with respect to the primal variables gives
$$
\wv = \sum_{i=1}^n \alpha_i y^{(i)} \xv^{(i)},
\qquad
\sum_{i=1}^n \alpha_i y^{(i)} = 0,
\qquad
C - \alpha_i - \gamma_i = 0.
$$
\item The equality constraint comes from differentiating with respect to the intercept $b$.
\item Since $\gamma_i \ge 0$, the dual variables satisfy the box constraints
$$
0 \le \alpha_i \le C.
$$
\item Key interpretation: the optimal weight vector is a linear combination of training points.
\end{framei}


\begin{framei}[fs=small]{ML Example: Resulting SVM Dual}
\item Eliminating the primal variables leaves the dual quadratic program
$$
\max_{\substack{0 \le \alpha_i \le C \\ i=1,\dots,n}}
\left[
\sum_{i=1}^n \alpha_i
-
\frac{1}{2}
\sum_{i=1}^n \sum_{j=1}^n
\alpha_i \alpha_j y^{(i)} y^{(j)} {\xv^{(i)}}^\top \xv^{(j)}
\right].
$$
\vspace*{-0.3cm}
$$
\text{s.t. } \qquad \sum_{i=1}^n \alpha_i y^{(i)} = 0.
$$
\item This is a \textbf{box-constrained QP plus one equality constraint} in the dual variables only.
\item Once the optimal multipliers are known, the primal separator follows from
$$
\wv^\ast = \sum_{i=1}^n \alpha_i^\ast y^{(i)} \xv^{(i)}.
$$
\item Only points with $\alpha_i^\ast > 0$ contribute to $\wv^\ast$; these are the \textbf{support vectors}.
\item The dual is especially attractive because it depends on the data only through dot products ${\xv^{(i)}}^\top \xv^{(j)}$, which later enables kernel methods.
\end{framei}


\begin{framev}{Take-Away}
\begin{enumerate}
\item Lagrangian relaxation converts constraints into multiplier variables.
\item The dual value is a lower bound:
$$
d^\ast = \max_{\alphav \ge \zero} \min_{\xv} H(\xv,\alphav) \le p^\ast.
$$
\item The primal uses the same objective in reversed order:
$$
p^\ast = \min_{\xv} \max_{\alphav \ge \zero} H(\xv,\alphav).
$$
\item For convex problems with Slater's condition, strong duality and the KKT conditions let us solve and certify optimality through the dual.
\end{enumerate}
\vfill
\oneliner{Dual methods are especially attractive when the Lagrangian lets us eliminate many primal variables and leaves a simpler problem in the multipliers.}
\end{framev}

\endlecture
\end{document}
